% PCT-SOURCE: pdf=038 print=026 kind=figure id=1-3
% Native reconstruction from the rendered source page.  The vertical hatch
% marks the two-particle continuum and the crossed hatch marks the
% three-particle continuum.
\begingroup
\renewcommand{\thefigure}{1-3}
\renewcommand{\figurename}{FIGURE}
\begin{figure}[t]
  \centering
  \begin{tikzpicture}[x=0.78cm,y=0.78cm,>=Stealth,
    line cap=round,line join=round]
    \begin{scope}[xshift=2.60cm]
    % The four straight spectral-cone rays.
    \draw[line width=.78pt] (-4.30,5.55) -- (0,0) -- (4.30,5.55);
    \draw[line width=.78pt] (-3.78,5.30) -- (0,0) -- (3.78,5.30);
    \draw[line width=.60pt] (-0.62,-0.62) -- (0,0) -- (0.62,-0.62);

    % Two-particle continuum: vertical hatching between the two-particle
    % threshold and the upper boundary.
    \begin{scope}
      \clip
        plot[domain=-3.55:3.55,samples=100]
          (\x,{3.02+0.18*\x*\x})
        -- (3.55,{4.22+0.12*3.55*3.55})
        plot[domain=3.55:-3.55,samples=100]
          (\x,{4.22+0.12*\x*\x}) -- cycle;
      \foreach \x in {-3.6,-3.35,...,3.6}
        \draw[line width=.38pt] (\x,2.85) -- (\x,5.95);
    \end{scope}

    % Three-particle continuum: crossed hatching above the upper boundary.
    \begin{scope}
      \clip
        plot[domain=-3.55:3.55,samples=100]
          (\x,{4.22+0.12*\x*\x})
        -- (3.55,6.05) -- (-3.55,6.05) -- cycle;
      \draw[step=.28cm,line width=.34pt] (-3.65,4.10) grid (3.65,6.10);
    \end{scope}

    % Curves bounding the mass shells and continua.
    \draw[line width=.85pt,domain=-3.55:3.55,samples=140]
      plot (\x,{1.48+0.25*\x*\x});
    \draw[line width=.85pt,domain=-3.55:3.55,samples=140]
      plot (\x,{3.02+0.18*\x*\x});
    \draw[line width=.85pt,domain=-3.55:3.55,samples=140]
      plot (\x,{4.22+0.12*\x*\x});

    % Axes, labels, and the vacuum point.
    \draw[->,line width=.60pt] (0,-0.25) -- (0,6.10)
      node[above right=1pt] {$p^0$};
    \draw[->,line width=.60pt] (0,0) -- (4.52,0)
      node[right=2pt] {$p$};
    \fill (0,0) circle (1.8pt);
    \node[below right=3pt] at (0,0) {$\Omega$};

    \draw[line width=.45pt] (0.63,1.60) -- (1.42,0.85);
    \node[right] at (1.50,0.80) {$-p^2=m^2$};

    \node[anchor=west] at (3.25,1.42) {one-particle};
    \node[anchor=west] at (3.25,1.02) {states $[m,0]$};
    \end{scope}

    % Legend for the two hatching patterns.
    \foreach \x in {-3.92,-3.70,-3.48,-3.26}
      \draw[line width=.55pt] (\x,3.35) -- (\x,3.78);
    \node[anchor=west] at (-4.08,3.02) {two-particle states};
    \node[anchor=west] at (-4.08,2.62)
      {$([m,0]\otimes[m,0])_{\mathrm{s}}$};
    \foreach \y in {1.70,1.92,2.14,2.36}
      \draw[line width=.55pt] (-3.92,\y) -- (-3.26,\y);
    \node[anchor=west] at (-4.08,1.38) {three-particle states};
    \node[anchor=west] at (-4.08,0.98)
      {$([m,0]\otimes[m,0]\otimes[m,0])_{\mathrm{s}}$};
  \end{tikzpicture}
  \caption{The spectrum of a theory of neutral scalar mesons of mass $m$
  without bound states.}
  \label{fig:1-3}
\end{figure}
\endgroup
